% ===================== CAPÍTULO 4 - ANÁLISE DE RESULTADOS =====================
\chapter{Análise de Resultados}\label{chap:analise}

\section{Cumprimento dos Objetivos}\label{sec:cumprimento}

De forma geral, os objetivos definidos para o estágio foram cumpridos, tendo sido necessário ajustar algumas prioridades na fase final do projeto para responder a desafios técnicos entretanto identificados. A avaliação do cumprimento dos objetivos organiza-se em torno das três macrofases de desenvolvimento que estruturaram o projeto.

O \textit{backend} - correspondente às Fases~1 a~3 do planeamento - cumpriu integralmente o objetivo central de integrar o fornecedor SeaRates à \textit{Tracking \gls{api}} existente. O mecanismo de orquestração (\textit{ResolveProviderService}) foi estendido para suportar o encaminhamento determinístico de pedidos ao novo fornecedor, sem qualquer alteração nos \textit{endpoints} públicos já expostos a clientes externos (\autoref{fig:orchestrator}).

A validação em ambiente simulado, com recurso ao Prism para \textit{mocking} da \gls{api} SeaRates, contribuiu para garantir a estabilidade da integração e a conformidade com os contratos de dados definidos por meio dos \textit{\glspl{dto}} implementados, cumprindo o objetivo de expandir a \gls{api} sem comprometer a retrocompatibilidade com os consumidores existentes.

A Fase~4 produziu uma interface \textit{frontend} que ultrapassou o propósito inicial de uma ferramenta de consulta de \textit{tracking}, consolidando-se como um painel administrativo completo em React (\anexoref{anx:solucao}). Foram implementadas e testadas operações \textit{\gls{crud}} dinâmicas baseadas em metadados para as várias entidades do sistema, o controlo granular de acessos através de enumerados, bem como a gestão detalhada de créditos e de integrações visuais. As funcionalidades implementadas permitiram disponibilizar uma plataforma utilizável tanto por clientes finais como por utilizadores administrativos.

As Fases~5 e~6, dedicadas às integrações e à estabilização, partiam de um plano que previa uma integração aprofundada com a aplicação corporativa i-Forwarder. Embora este objetivo não tenha sido atingido na sua totalidade, a reorientação estratégica das prioridades mostrou-se acertada. O foco foi redirecionado para as necessidades críticas da plataforma, resultando numa \gls{api} estabilizada, com tratamento de exceções específico para a SeaRates e mensagens de movimento normalizadas.

A migração do sistema de autenticação para Keycloak e a correção de anomalias na renderização geográfica garantiram que a solução final cumprisse os requisitos de segurança, fiabilidade e usabilidade exigidos para adoção interna.

Do ponto de vista quantitativo, o estágio totalizou 416 horas de trabalho efetivo ao longo de 16 semanas. Ao fornecedor ShipsGo~\cite{shipsgo}, já existente antes do estágio, foram acrescentados dois novos \textit{providers}: o SeaRates~\cite{searates}, integrado pelo autor, e o OpenSyllo (\gls{api} privada, sem documentação pública acessível), integrado em paralelo pelo colega de equipa Sandro Cunha. A comunicação com os fornecedores externos é uniformizada por meio de três interfaces base (\textit{ITrackingCoreProvider}, \textit{ICarrierProvider} e \textit{IAirlineProvider}).

A \gls{api} expõe 37 \textit{controllers}, organizados em 10 domínios funcionais, e é suportada por 26 entidades na \gls{bd} do SQL Server. No \textit{frontend}, foram implementadas operações \textit{\gls{crud}} para 14 entidades, com suporte para três idiomas (português, inglês e espanhol). A qualidade do código foi assegurada por aproximadamente 55 ficheiros de testes unitários e de integração (xUnit~\cite{xunit}, Moq~\cite{moq} e FluentAssertions~\cite{fluentassertions}), com uma \textit{test suite} que cobre os principais fluxos do orquestrador de \textit{providers} e dos serviços de \textit{tracking}.

A \autoref{tab:comparacao-horas} sintetiza os desvios entre as horas planeadas e as horas efetivamente executadas em cada uma das seis fases do estágio.

\begin{table}[H]
\centering
\caption{Comparação entre horas planeadas e executadas por fase}
\label{tab:comparacao-horas}
\small
\begin{tabular}{|p{0.42\textwidth}|c|c|c|}
\hline
\textbf{Fase} & \textbf{Planeadas (h)} & \textbf{Reais (h)} & \textbf{Desvio} \\
\hline
Fase 1 - Formação Técnica na \textit{stack} Devlop & 32 & 32 & 0\,\% \\
\hline
Fase 2 - Análise da \gls{api} existente e arquitetura atual & 40 & 24 & $-$40\,\% \\
\hline
Fase 3 - Integração do novo fornecedor SeaRates & 160 & 72 & $-$55\,\% \\
\hline
Fase 4 - Desenvolvimento do \textit{frontend} da Plataforma de Tracking & 320 & 168 & $-$48\,\% \\
\hline
Fase 5 - Integração com Aplicação Interna (i-Forwarder) & 80 & 68 & $-$15\,\% \\
\hline
Fase 6 - Testes finais, documentação e conclusão & 40 & 52 & +30\,\% \\
\hline
\textbf{Total} & \textbf{672} & \textbf{416} & \textbf{$-$38\,\%} \\
\hline
\end{tabular}
\end{table}

\section{Discussão}

A execução do plano de trabalhos decorreu, no essencial, em conformidade com a calendarização prevista no \gls{pit} (\anexoref{anx:pit}), registando-se, contudo, desvios na fase final que merecem reflexão. O principal desvio face ao planeamento inicial consistiu na redução do âmbito da integração com o sistema i-Forwarder, em favor de um investimento mais aprofundado na estabilização da plataforma e no reforço da sua segurança.

Esta decisão, tomada em articulação com o orientador da empresa, foi considerada a mais adequada tendo em conta o estado de desenvolvimento do produto.

Entre as dificuldades técnicas encontradas, destaca-se a resolução do problema do antimeridiano na renderização geográfica de trajetos marítimos, cuja descrição técnica foi apresentada no \autoref{chap:descricao}. A migração do sistema de autenticação para Keycloak também representou um desafio relevante, exigindo adaptações profundas na lógica de autorização e no fluxo de autenticação do \textit{frontend} e do \textit{backend}. Em ambos os casos, os obstáculos foram superados com sucesso, o que contribuiu para a robustez da solução final.

Os objetivos gerais do estágio foram cumpridos: a integração do novo fornecedor SeaRates ficou plenamente operacional e a plataforma de \textit{tracking} dispõe de uma interface funcional e segura. Os desvios em relação ao planeamento inicial foram geridos de forma adaptativa, com decisões orientadas pela qualidade e estabilidade da solução.

\section{Tecnologias e Ferramentas}\label{sec:tecnologias}

O conjunto de tecnologias e ferramentas que suportaram o desenvolvimento do projeto encontra-se sistematizado na \autoref{tab:tecnologias}, organizado por categoria e respetivo âmbito de aplicação.

\begin{table}[H]
\centering
\caption{Tecnologias e ferramentas utilizadas no estágio}
\label{tab:tecnologias}
\small
\begin{tabular}{@{}p{0.20\textwidth} p{0.35\textwidth} p{0.37\textwidth}@{}}
\toprule
\textbf{Categoria} & \textbf{Tecnologia / Ferramenta} & \textbf{Âmbito} \\
\midrule
\multirow{2}{=}{Linguagens}
  & C\# & \textit{backend}: \gls{api}, serviços de \textit{tracking}, lógica aplicacional \\
  & JavaScript (React) & \textit{frontend}: interface da plataforma \\
\midrule
\multirow{3}{=}{Frameworks}
  & .NET \cite{dotnet} & Base da \gls{api} e lógica de negócio \\
  & \gls{efcore} \cite{efcore} & Acesso a dados (SQL Server) \\
  & React \cite{react} & Construção da interface de \textit{tracking} \\
\midrule
\multirow{2}{=}{Testes}
  & xUnit, Moq, FluentAssertions \cite{xunit,moq,fluentassertions} & Testes unitários e de integração \\
  & Prism \cite{prism} & Simulação da \gls{api} SeaRates \\
\midrule
\multirow{2}{=}{Ambientes de desenvolvimento}
  & JetBrains Rider \cite{rider} & \textit{backend} (.NET) \\
  & JetBrains WebStorm \cite{webstorm} & \textit{frontend} (React) \\
\midrule
\multirow{4}{=}{Infraestrutura}
  & Microsoft SQL Server \cite{sqlserver} & Base de dados \\
  & Git / Azure DevOps \cite{azuredevops} & Controlo de versões e colaboração \\
  & Keycloak \cite{keycloak} & Autenticação e gestão de identidades \\
  & Four Js Genero (\gls{4gl}) \cite{fourjs} & Integração com i-Forwarder \\
\midrule
\multirow{4}{=}{Ferramentas de apoio}
  & Swagger \cite{swagger} & Teste e exploração da \gls{api} \\
  & GitHub Copilot \cite{copilot} / OpenCode Desktop~\cite{opencode} & Apoio à geração de código (ver \ref{sec:ia}) \\
  & Microsoft Teams & Comunicação e acompanhamento \\
  & Microsoft Word & Relatórios e documentação \\
\bottomrule
\end{tabular}
\end{table}

A \autoref{tab:tecnologias} sistematiza o conjunto de tecnologias e ferramentas utilizadas. Destacam-se, pelo seu papel transversal no desenvolvimento, as ferramentas de inteligência artificial, abordadas em maior detalhe de seguida.

\subsection{Utilização de Ferramentas de Inteligência Artificial}\label{sec:ia}

Ao longo do estágio, recorreu-se a diversas ferramentas de inteligência artificial como apoio ao desenvolvimento. O GitHub Copilot~\cite{copilot}, integrado nos ambientes JetBrains Rider e WebStorm, foi utilizado de forma transversal para agilizar a escrita de código repetitivo, a geração de estruturas de dados e a sugestão de implementações. Estas ferramentas revelaram-se particularmente úteis na integração com a SeaRates, acelerando a criação dos \textit{\glspl{dto}} de pedido e resposta, dos modelos de domínio e dos testes unitários para os novos serviços de \textit{tracking}. Para questões mais abrangentes de arquitetura, depuração de erros complexos e obtenção de trechos de código específicos, recorreu-se também ao ChatGPT~\cite{chatgpt} e ao Gemini~\cite{gemini}.

Nas últimas semanas de cada ciclo de trabalho, quando os limites de utilização do GitHub Copilot eram atingidos, o Gemini Code Assist serviu como ferramenta de recurso para tarefas pontuais de código. Contudo, nas duas últimas semanas do estágio, a imposição de restrições mais severas aos \textit{tokens} do Copilot coincidiu com a fase final de desenvolvimento do \textit{frontend} - onde o volume de trabalho de implementação de componentes React, a correção da anomalia do antimeridiano e a migração para Keycloak exigiam um ritmo de produção elevado.

Para fazer face a esta contingência e manter a produtividade sem interrupções, adotou-se a versão desktop do OpenCode~\cite{opencode} com o modelo DeepSeek~\cite{deepseek} como alternativa principal. Esta combinação constituiu uma alternativa eficaz, permitindo dar continuidade ao trabalho e assegurar o cumprimento dos prazos estabelecidos para a conclusão do estágio. Adicionalmente, estas ferramentas foram utilizadas como apoio na configuração e na estruturação do presente relatório em LaTeX.

